\clearpage
\VerdierPhysicalPageAnchor{0170}%
\begingroup\small\setlength{\parskip}{0pt}\setlength{\emergencystretch}{2.5em}
\VerdierSemanticAnchor{ch4-4-2-7-second-diagram}
Applying the same observation to the diagram
\[
\begin{tikzcd}[column sep=4.6em,row sep=2.8em]
&H^n(f_2f_1)\arrow[d]\arrow[dr]&\\
H^n(f_1)\arrow[ur]\arrow[r]&
H^n(f_3f_2f_1)\arrow[r]&H^n(f_3f_2),
\end{tikzcd}
\]
gives an isomorphism
\VerdierSemanticAnchor{ch4-4-2-7-2}
\[
E\!\begin{bmatrix}n\\ f_1,f_2,f_3\end{bmatrix}
\xrightarrow{\ \sim\ }
\operatorname{Im}\bigl(H^n(f_2f_1)\longrightarrow H^n(f_3f_2f_1)\bigr)
\Big/
\operatorname{Im}\bigl(H^n(f_1)\longrightarrow H^n(f_3f_2f_1)\bigr).
\tag{4.2.7.2}
\]

\VerdierSemanticAnchor{ch4-4-3}
\medskip\noindent\textbf{4.3. The terms
$E\!\begin{bmatrix}n\\ -\infty,p,q,+\infty\end{bmatrix}$.}

\VerdierSemanticAnchor{ch4-4-3-1}
\medskip\noindent\textbf{4.3.1.}
Let $\widetilde{\mathbb Z}$ be the ordered set $\mathbb Z$ completed by an
initial element $-\infty$ and a final element $+\infty$. By abuse of
notation, we shall also use $\widetilde{\mathbb Z}$ for the category
associated with the ordered set $\widetilde{\mathbb Z}$. A morphism of
$\widetilde{\mathbb Z}$ is therefore a pair $p\leq q$, where $p$ and $q$
are elements of $\widetilde{\mathbb Z}$. An ordered sequence
$p_1\leq p_2\leq\cdots\leq p_n$ defines an object of
$\mathcal{F}\ell_{n-1}(\widetilde{\mathbb Z})$, denoted
$(p_1,p_2,\ldots,p_n)$. A morphism of $\widetilde{\mathbb Z}$ is called
\textit{finite} if its source and target belong to $\mathbb Z$. An object
of $\mathcal{F}\ell_n(\widetilde{\mathbb Z})$ is called \textit{finite} if
its constituent morphisms are finite.

\VerdierSemanticAnchor{ch4-4-3-2}
\medskip\noindent\textbf{4.3.2.}
\textit{In this work we shall use only spectral objects of type
$\widetilde{\mathbb Z}$; henceforth these will simply be called spectral
objects.} This recovers the notion of [1], Chapter XV, §7. Let
$(H^n(p,q),\delta,\ p\leq q\in\widetilde{\mathbb Z})$ be a spectral object
with values in an abelian category $\mathcal A$. Set
\VerdierSemanticAnchor{ch4-4-3-2-1}
\[
F^pH^n(-\infty,+\infty)
=\operatorname{Im}\bigl(H^n(-\infty,p)\longrightarrow
H^n(-\infty,+\infty)\bigr).
\tag{4.3.2.1}
\]
We have $F^{+\infty}H^n(-\infty,+\infty)=H^n(-\infty,+\infty)$ and
$F^{-\infty}H^n(-\infty,+\infty)=0$. The
$F^pH^n(-\infty,+\infty)$ form an increasing filtration of type
$\widetilde{\mathbb Z}$ on $H^n(-\infty,+\infty)$. Moreover, for every pair
$p\leq q$, there is an isomorphism (4.2.7.2)
\VerdierSemanticAnchor{ch4-4-3-2-filtration-quotient}
\[
E\!\begin{bmatrix}n\\ -\infty,p,q,+\infty\end{bmatrix}
\simeq
F^qH^n(-\infty,+\infty)\big/F^pH^n(-\infty,+\infty).
\]
\endgroup
